\documentclass[output=paper]{langscibook}
\ChapterDOI{10.5281/zenodo.22880250}
\title{\textit{O kara roge} – Hibiscus fibres for making strings}
\author{Joyce Maion}
\abstract{}
\begin{document}
\maketitle\label{chap:12}
\begin{figure}[H]
    \captionsetup[subfigure]{oneside,margin={0cm,0cm}}
    \subfloat[]{\includegraphics[trim=0 0 0 0,clip,width=.43\linewidth]{figures/12_1.jpg}\label{fig:12_1}}
    \subfloat[]{\includegraphics[trim=0 0 0 50,clip,width=.28\linewidth]{figures/12_2.jpg}\label{fig:12_2}}
    \subfloat[]{\includegraphics[trim=0 160 0 0,clip,width=.29\linewidth]{figures/12_3.jpg}\label{fig:12_3}}
    \caption{(a) \textit{O roge} -- Coastal hibiscus, (b) \textit{O taporohi roge} -- A hibiscus flower, (c) \textit{Kara roge} -- Strings made of the hibiscus fibres}
\end{figure}
\begin{xltabular}{\textwidth}{lXX}
    (\ref{ex:text8-1})&\textit{Bean rake nom to kara roge,}&If you want a string made of hibiscus\footnote{\textit{Hibiscus tiliaceus}} fibres,\\
    (\ref{ex:text8-2})&\textit{ean toro nao momohu tea mosi to vehu roge vai o kobus.}&you must first go and cut a moist hibiscus twig.\\
    (\ref{ex:text8-3})&\textit{O vehu roge vai na tahihi vakahu nana.}&The fibres of hibiscus twig are quickly peeled off.\\
    (\ref{ex:text8-4})&\textit{Ean repaa hihi a kapa nae.}&Then you peel its cover off.\\
    (\ref{ex:text8-5})&\textit{Hihi kahii bono naono nae.}&You peel it off its wood.\\
    (\ref{ex:text8-6})&\textit{O kapa vai to hihi noman, eve to dao raara bono kara roge.}&The cover that you peel off, we call hibiscus fibre.\\
    (\ref{ex:text8-7})&\textit{Benam hiki to urisi mene varoto bona vateen,}&When we don't have a wild banana sheath to strap up the backpack,\\
    (\ref{ex:text8-8})&\textit{enam repaa varoto nio kara roge.}&we strap it up with a strap made of hibiscus fibres.\\
    (\ref{ex:text8-9})&\textit{Meha maa vuri, eori repaa paku bari bono sikioti vaa teo hagi.}&Sometimes they make the skirt for a dance from them.\\
    (\ref{ex:text8-10})&\textit{Hihi rori bona kapa roge,}&They peel the hibiscus covers off,\\
    (\ref{ex:text8-11})&\textit{eori repaa taono ni bari tahii.}&then they put them in the sea, weighing them down with stones.\\
    (\ref{ex:text8-12})&\textit{Beo kapa roge vai o pereperee koava,}&When a green hibiscus cover rots,\\
    (\ref{ex:text8-13})&\textit{eori repaa gono takin kahi bari bona tahii.}&they get it out of the saltwater.\\
    (\ref{ex:text8-14})&\textit{Murinae eori repaa garuhu bari,}&After that they wash it,\\
    (\ref{ex:text8-15})&\textit{eori repaa varoava bari.}&and put it in the sun.\\
    (\ref{ex:text8-16})&\textit{Beve raka, eori repaa sirisiri varutarutaa bari.}&When it is dry, they tear it into fine strips.\\
    (\ref{ex:text8-17})&\textit{Eori repaa paku bari bono garasi sikioti mene hagi teori.}&Then they use them for making their grass skirt for dancing.\\
\end{xltabular}
\end{document}